% --- Άσκηση 38791 (38791.tex) ---
\Askhsh{38791}{4}
\begin{ylh}{2}
    \item 6.2. Γραφική Παράσταση Συνάρτησης
    \item 6.3. Η Συνάρτηση $f(x)=ax+\beta$
\end{ylh}

Ένας μαραγκός έχει μια μικρή επιχείρηση που φτιάχνει και πουλάει ξύλινα τραπέζια. Επιθυμώντας να αυξήσει τα κέρδη του, αποφάσισε αρχικά να μελετήσει την παραγωγή και τις πωλήσεις του.

Υπολόγισε ότι, όταν πωλούνται $x$ τραπέζια, η τιμή πώλησης σε ευρώ ενός τραπεζιού είναι $p(x) = 120 - x$, με $0 \leq x \leq 100$.

\begin{alist}
\item Ποιος είναι ο τύπος που εκφράζει τα έσοδα της επιχείρησης από την πώληση $x$ τραπεζιών;
\monades{3}
Στη συνέχεια βρήκε επίσης ότι το μηνιαίο κόστος παραγωγής $K(x)$ σε ευρώ, όταν παράγει $x$ τραπέζια, προκύπτει από τη σχέση: $K(x) = 20x + 475$.

\item Να δώσετε μια ερμηνεία για το τι εκφράζει ο αριθμός $20$ και τι ο αριθμός $475$ στον τύπο του $K(x)$.
\monades{4}
\item Παρακάτω δίνονται οι γραφικές παραστάσεις των εσόδων $E(x)$ και του μηνιαίου κόστους παραγωγής $K(x)$ ως συνάρτηση της πώλησης $x$ τραπεζιών.

\begin{center}
\begin{tikzpicture}
\begin{axis}[
    axis lines = middle,
    xlabel = {τραπέζια $(x)$},
    ylabel = {\euro},
    xmin = -10, xmax = 115,
    ymin = -200, ymax = 3800,
    xtick = {10,20,30,40,50,60,70,80,90,100,110},
    ytick = {500,1000,1500,2000,2500,3000,3500},
    grid = both,
    grid style = {gray!30},
    samples = 100,
    width = 10cm, height = 7cm,
    tick label style = {font=\scriptsize},
    every axis x label/.style = {at={(current axis.right of origin)}, anchor=north west},
    every axis y label/.style = {at={(current axis.above origin)}, anchor=south},
]
